\documentclass[../../../include/open-logic-section]{subfiles}

\begin{document}

\olfileid{sth}{ord-arithmetic}{mult}

\olsection{ضربِ ترتیبی}

اکنون به ضرب ترتیبی می‌پردازیم و آن را تقریباً همانند جمع ترتیبی
بررسی می‌کنیم. فرض کنید می‌خواهیم $\alpha$ را در~$\beta$ ضرب کنیم.
برای این کار می‌توانید شبکه‌ای مستطیلی با پهنای $\alpha$ و ارتفاع
$\beta$ تصور کنید؛ حاصل‌ضربِ $\alpha$ و $\beta$ با پیمودن هر سطر و
سپس رفتن به سطر بعدی به دست می‌آید، \ldots، تا آنکه سراسر شبکه را
پیموده باشید. به بیان دیگر، حاصل‌ضربِ $\alpha$ و $\beta$ از آنجا
پدید می‌آید که \emph{هر} عضوِ $\beta$ را با نسخه‌ای از $\alpha$
جایگزین کنیم.

برای صورت‌بندی دقیق، کافی است ترتیب واژه‌نامه‌ایِ معکوس را روی
ضرب دکارتیِ $\alpha$ و $\beta$ به کار ببریم:

\begin{defn}
برای هر دو عدد ترتیبیِ $\alpha, \beta$، حاصل‌ضرب آن‌ها برابر است با
$\alpha \ordtimes \beta =
\ordtype{\alpha \times \beta, \rlexless}$.
\end{defn}
\noindent
باز باید تأیید کنیم که این تعریف خوش‌صورت است:

\begin{lem}\ollabel{ordtimeslessiswo}
$\tuple{\alpha \times \beta, \rlexless}$ برای هر دو عدد ترتیبیِ
$\alpha$ و $\beta$ یک خوش‌ترتیب است.
\end{lem}

\begin{proof}
دقیقاً مانند \olref[add]{ordsumlessiswo}.
\end{proof}
\noindent
اثبات رفتار بازگشتی زیر نیز دشوار نیست:

\begin{lem}\ollabel{ordtimesrecursion}
برای هر دو عدد ترتیبیِ $\alpha, \beta$:
\begin{align*}
	\alpha \ordtimes 0 &= 0\\
	\alpha \ordtimes (\beta \ordplus 1) &=
		(\alpha \ordtimes \beta) \ordplus \alpha\\
	\alpha  \ordtimes \beta &=
		\supstrict_{\delta < \beta}(\alpha \ordtimes \delta) &&
		\text{هرگاه $\beta$ عدد ترتیبیِ حدی باشد}.
\end{align*}
\end{lem}

\begin{proof}
به‌عنوان تمرین واگذار می‌شود.
\end{proof}

در واقع، درست مانند حالت جمع، می‌توانستیم ضرب ترتیبی را با این
معادلات بازگشتی تعریف کنیم، نه با یک تعریف مستقیم. باز هم مانند جمع،
برخی رفتارها آشنا هستند:

\begin{lem}\ollabel{ordinalmultiplicationisnice}
اگر $\alpha, \beta, \gamma$ اعداد ترتیبی باشند، آن‌گاه:
\begin{enumerate}
	\item\ollabel{ordtimes1} اگر $\alpha \neq 0$ و
	$\beta < \gamma$، آن‌گاه
	$\alpha \ordtimes \beta < \alpha \ordtimes \gamma$؛
	\item\ollabel{ordtimes2} اگر $\alpha \neq 0$ و
	$\alpha \ordtimes \beta = \alpha\ordtimes\gamma$، آن‌گاه
	$\beta = \gamma$؛
	\item\ollabel{ordtimes3}
	$\alpha \ordtimes (\beta \ordtimes \gamma) =
	(\alpha \ordtimes \beta) \ordtimes \gamma$؛
	\item\ollabel{ordtimes4} اگر $\alpha \leq \beta$، آن‌گاه
	$\alpha \ordtimes \gamma \leq \beta \ordtimes \gamma$؛
	\item\ollabel{ordtimes5}
	$\alpha \ordtimes (\beta \ordplus \gamma) =
	(\alpha \ordtimes \beta )\ordplus
	(\alpha\ordtimes \gamma)$.
\end{enumerate}
\end{lem}

\begin{proof}
به‌عنوان تمرین واگذار می‌شود.
\end{proof}

می‌توانید نتایج دیگری را به دلخواه اثبات کنید (یا در منابع
بیابید). اما با توجه به
\olref[ord-arithmetic][add]{ordsumnotcommute}، نتیجهٔ زیر نباید
غافلگیرکننده باشد:

\begin{prop}
ضرب ترتیبی جابجایی‌پذیر نیست:
$2 \ordtimes \omega = \omega < \omega \ordtimes 2$
\end{prop}

\begin{proof}
$2 \ordtimes \omega = \supstrict_{n < \omega}(2\ordtimes n) =
\omega \in \supstrict_{n < \omega}(\omega \ordplus n) =
\omega \ordplus \omega = \omega \ordtimes 2$.
\end{proof}
\noindent
دلیل شهودی باز هم کاملاً سرراست است. برای محاسبهٔ
$2 \ordtimes \omega$، هر عدد طبیعی را با دو شیء جایگزین می‌کنید.
اگر صرفاً همهٔ عددهای «نیمه» را میان اعداد طبیعی درج کنید نیز همان
نوع ترتیب را به دست می‌آورید؛ یعنی ترتیب طبیعی روی
$\Setabs{\nicefrac{n}{2}}{n \in \omega}$ را در نظر بگیرید. با این
بیان، آشکار است که نوع ترتیب با نوع ترتیبِ خودِ $\omega$ یکسان است.
اما برای محاسبهٔ $\omega \ordtimes 2$، دو نسخه از $\omega$ را
یکی پس از دیگری می‌گذارید.

\begin{prob}
\olref[sth][ord-arithmetic][mult]{ordtimeslessiswo}،
\olref[sth][ord-arithmetic][mult]{ordtimesrecursion}،
و
\olref[sth][ord-arithmetic][mult]{ordinalmultiplicationisnice}
را اثبات کنید.
\end{prob}

\end{document}
